\subsection{Projective properly elliptic components in zero-Euler fillings}

\begin{statement}{C.2.UN.18}{Proved}{Opposite logarithmic data produce a projective essential component}
Put $q=e^{-2\pi}$ and $E=\C^*/q^{\Z}$. Choose distinct
points $p_+,p_-\in E$ whose difference is torsion.
For every $m\ge2$ the fixed punctured family
$(E\times E)\times\Delta^*$ has a smooth proper filling
with full central fibre
\[
                         P_m+mQ_++mQ_-.                 \tag{C.PE.1}
\]
Here $Q_+$ and $Q_-$ are disjoint primary Hopf surfaces,
$P_m$ is a smooth projective surface of Kodaira dimension
one, and $J_\pm=P_m\cap Q_\pm$ are disjoint smooth
elliptic curves. The support is a simple normal crossing
surface of Euler number zero. The map $P_m\to E$ is an
elliptic fibration with exactly the two multiple fibres
$mJ_+$ and $mJ_-$, and
\[
 \begin{gathered}
 K_{P_m}\simeq\OO_{P_m}((m-1)(J_++J_-)),\qquad
 H_1(P_m,\Z)\simeq\Z^4,\\
 h^0(P_m,K_{P_m}^{\otimes r})=
 \begin{cases}1,&d_r=0,\\d_r,&d_r>0,\end{cases}
 \qquad d_r=2\left\lfloor\frac{r(m-1)}m\right\rfloor.
 \end{gathered}                                         \tag{C.PE.2}
\]
The surface $P_m$ admits a finite \'etale cover by a
product of two projective curves, one elliptic and the
other of genus at least two. All three components contain
no rational curve and survive birationally in every
bimeromorphic comparison between smooth fillings. The
multiplicity-one component $P_m$ also survives birationally
after every finite base change and resolution. The
fillings for different $m$ are not bimeromorphic over the
disc.

The construction extends to a smooth compact threefold
$Y_m\to\PP^1$ with just this special fibre. All integral
cohomology groups below are free. Writing $S_m$ for its
reduced special support, their ranks and coherent
cohomology dimensions are
\[
 \begin{array}{c|c|c}
 &\operatorname{rank}H^k(-,\Z)&h^k(\OO)\\ \hline
 P_m&(1,4,6,4,1)&(1,2,1)\\
 Q_\pm&(1,1,0,1,1)&(1,1,0)\\
 S_m&(1,3,6,7,3)&(1,2,1)\\
 Y_m&(1,3,7,10,7,3,1)&(1,2,1,0).
 \end{array}                                           \tag{C.PE.3}
\]
\end{statement}
\begin{proof}
Let $V$ be the blowup of $E\times\Delta$ at
$(p_+,0)$ and $(p_-,0)$, and let $R_+,R_-$ be its
disjoint exceptional curves. Put
\[
 D=R_+-R_-,\qquad L=\OO_V(D),\qquad
 \operatorname{div}(\sigma)=D,                         \tag{C.PE.4}
\]
where $\sigma$ is the canonical meromorphic section of
$L$. We define a smooth covering space $W_m$ using two
actual charts. On $V\setminus R_-$ use
\[
 \sigma=a^m\ell,\qquad \ell\in L^\times,\quad a\in\C;
                                                               \tag{C.PE.5}
\]
on $V\setminus R_+$ use
\[
 \sigma^{-1}=(a')^m\ell',\qquad
 \ell'\in(L^{-1})^\times,\quad a'\in\C.                 \tag{C.PE.6}
\]
On the common open set both $a,a'$ are nonzero, and the
gluing is $a'=a^{-1}$, $\ell'=\ell^{-1}$. Thus $a$
is globally a $\PP^1$-valued coordinate. The generator
of the free $\Z$-action is
\[
 (\ell,a)\longmapsto(q^m\ell,q^{-1}a),\qquad
 (\ell',a')\longmapsto(q^{-m}\ell',qa').                 \tag{C.PE.7}
\]
Near $R_+$, a frame of $L$ makes (C.PE.5) the equation
$x=a^m b$, $b\ne0$. Near $R_-$, a frame of $L^{-1}$
makes (C.PE.6) the equation $x=(a')^m b'$, $b'\ne0$.
Each eliminates $x$ and proves smoothness. A Hermitian
norm of the globally defined nonzero vector $\ell$ is
multiplied by $q^m$ in (C.PE.7); it proves proper
discontinuity and the absence of stabilizers. The
quotient $U_m=W_m/\Z$ is therefore a smooth threefold.

The local proper maps are explicit. In the positive
chart write $x=z^m$, $b=w^m$, $a=z/w$ and quotient
by $(z,w)\mapsto(\zeta z,\zeta w)$. The generator
(C.PE.7) lifts to $w\mapsto qw$. In the negative
chart use $x=z^m$, $b'=(w')^m$, $a'=z/w'$ and the
same diagonal finite action; the generator lifts to
$w'\mapsto q^{-1}w'$. In each case the local map to
$V$ is a free finite quotient of
\[
 \Delta_z\times E\times\Delta_y
   \longrightarrow\Delta_x\times\Delta_y,\qquad
 (z,[w],y)\longmapsto(z^m,y).                           \tag{C.PE.8}
\]
Multiplication by a nontrivial $m$th root of unity is a
nontrivial translation of $E$, so the finite action is
free. The map before this quotient is proper; hence
the resulting $\rho:U_m\to V$ is proper. Away from
$R_+\cup R_-$ the actual product identification is
\[
 (v,[c])\longmapsto[(v,\ell=\sigma(v)c^m,a=c^{-1})].     \tag{C.PE.9}
\]
It gives exactly the stated fixed punctured product.

Write $B\simeq E$ for the strict central curve of $V$.
The base function has divisor $B+R_++R_-$. At the two
intersections it has equation $t=xy$, so its pullbacks
are respectively
\[
                   t=a^mby,\qquad t=(a')^m b'y.         \tag{C.PE.10}
\]
These are transverse products with multiplicities
$m,1$. They prove (C.PE.1), smoothness of $P_m$ and
the stated intersections. Over $R_+$ the component is
$L|_{R_+}^{\times}/q^{m\Z}$; over $R_-$ it is
$(L^{-1}|_{R_-})^{\times}/q^{-m\Z}$. Both line bundles
are $\OO_{\PP^1}(-1)$. Sending a nonzero tautological
vector to the same vector in $\C^2\setminus0$ identifies
both components with
$(\C^2\setminus0)/q^{m\Z}$, a primary Hopf surface.
The two reduced multiple fibres $J_\pm$ are each
$\C^*/q^{m\Z}$.

We construct the finite projective cover of $P_m$.
The restriction of $L$ to $B$ is
$L_B=\OO_E(p_+-p_-)$, a torsion line. Choose a torsion
line $N$ with $N^{\otimes m}\simeq L_B$; multiplication
by $m$ is onto on the torsion subgroup of $\Pic^0(E)$.
Normalize the cyclic cover described in the projective
bundle $\PP(\OO_E\oplus N)$ by
\[
                             y^m=\sigma|_B.              \tag{C.PE.11}
\]
Here $y$ is the tautological meromorphic $N$-coordinate.
At $p_+$ this is $y^m=z$, and at $p_-$ its reciprocal
is $(y^{-1})^m=z$, after local frame choices. The
normalization $C$ is a connected smooth projective
curve: the order-one zero of $\sigma$ prevents the
cyclic equation from splitting, and the local maps
are the stated fully ramified ones. Riemann--Hurwitz
gives $2g(C)-2=2(m-1)$, so $g(C)=m$.

Pull back a finite \'etale cover of $E$ trivializing
$N$, and choose a connected component of its fibre
product with $C$. This gives a finite \'etale cover
$C'\to C$ and a map $\tau:C'\to E$ such that
$\tau^*N$ has a nowhere-zero frame $n$. Write
$y=y_0n$. There is a holomorphic map
\[
 \begin{split}
 \Phi:C'\times E&\longrightarrow P_m,\\
 (z,[u])&\longmapsto
       [\ell=u^m n(z)^m,\ a=y_0(z)/u].
 \end{split}                                             \tag{C.PE.12}
\]
At a pole of $y_0$ this formula is interpreted in
(C.PE.6), so $a'=u/y_0$ and $\ell'=u^{-m}n^{-m}$.
It is well defined under $u\mapsto qu$ by (C.PE.7).
It is onto and has finite fibres: away from the two
branch values its restriction to each elliptic fibre
is an isomorphism; at a branch value it is the finite
map $[u]\mapsto[u^m]$ from $E$ to
$\C^*/q^{m\Z}$. Thus it is finite and proper.

In fact $\Phi$ is \'etale everywhere. Local nonvanishing
frame factors are absorbed into the coordinate $u$.
At a zero branch choose the root parameter $\zeta$;
in the actual surface coordinates
$(a,b)$ its map is
\[
             (\zeta,u)\longmapsto(\zeta/u,u^m),
             \qquad \det d\Phi=m u^{m-2}\ne0.           \tag{C.PE.13}
\]
At a pole use $y_0=\zeta^{-1}$ and the map
$(\zeta,u)\mapsto(a'=\zeta u,b'=u^{-m})$, whose
Jacobian is $-m u^{-m}\ne0$. Away from those branches
the assertion follows from the \'etale base map and
the fibre isomorphism. This proves the claimed finite
\'etale product cover. Its source is projective, and
a finite holomorphic image of a projective surface is
Moishezon; a smooth compact Moishezon surface is
projective. The comparison can also be made by taking
the finite \'etale Galois closure of (C.PE.12), a
projective surface, and its finite algebraic quotient.

Let $\omega$ be a nowhere-zero one-form on $E$.
Off the multiple fibres, the form
$p^*\omega\wedge d\log\ell$ is well defined: a change
of frame adds a base one-form to $d\log\ell$, whose
wedge with $p^*\omega$ is zero. Near $J_+$, write
$\omega=h(x)dx$ and $x=a^m b$; then this two-form is
\[
                m h(a^m b)a^{m-1}\,da\wedge db.         \tag{C.PE.14}
\]
Near $J_-$ the dual formula has a minus sign and the
same vanishing order $m-1$. It has no further zero.
This proves the canonical line in (C.PE.2). Constancy
on compact elliptic fibres and the invariant local
root charts give the exact direct image
\[
 p_*\OO_{P_m}(k_+J_++k_-J_-)
 =\OO_E\!\left(
  \left\lfloor\frac{k_+}m\right\rfloor p_+
 +\left\lfloor\frac{k_-}m\right\rfloor p_-\right).         \tag{C.PE.15}
\]
A permitted Laurent term in a root parameter descends
precisely when its order is divisible by $m$.
Substitute $k_+=k_-=r(m-1)$ and use Riemann--Roch
on the elliptic curve. This proves the plurigenera
and $\kappa(P_m)=1$.

We give the integral comparisons in (C.PE.3). The
real action
$(\ell,a)\mapsto(\lambda^m\ell,\lambda^{-1}a)$,
$\lambda>0$, preserves the equations and translates
$\log\|\ell\|$. Since $q$ is positive real, its
quotient gives actual differentiable products
\[
                   P_m\simeq M_3\times S^1,\qquad
                   Y_m\simeq M_5\times S^1,             \tag{C.PE.16}
\]
where the $M$'s are the level sets $\|\ell\|=1$ in
the corresponding covering spaces. For $Y_m$ use
the construction over
$V=\operatorname{Bl}_{(p_+,0),(p_-,0)}(E\times\PP^1)$.
The local proper maps (C.PE.8) make this quotient
compact and smooth. Off $0$ the product comparison
(C.PE.9) applies, including at infinity.

For $M_3$, remove small discs around $p_+,p_-$.
Over their complement the angular coordinate of
$a$ trivializes its circle fibration. Let $h$ be
the positive angular $a$-loop and let $\delta_\pm$
be positive base boundary loops with $a$ fixed.
In the positive solid torus the phases satisfy
$\delta_+=m\arg(a)+\arg(b)$; in the negative one
they satisfy $\delta_-=m\arg(a')+\arg(b')$ and
$\arg(a)=-\arg(a')$. The two meridians are
\[
                         m\delta_++h,\qquad
                         m\delta_- -h.                 \tag{C.PE.17}
\]
In the abelianization of the punctured elliptic base,
$\delta_++\delta_-=0$. These relations give
$h=-m\delta_+$ and $\delta_-=-\delta_+$, leaving
$H_1(M_3,\Z)=\Z^3$, generated by the two
elliptic-base loops and $\delta_+$. There is no
torsion. Poincar\'e duality on this closed oriented
three-manifold and (C.PE.16) give all the asserted
integral groups of $P_m$.

For the compact five-manifold put
$Q=(E\times\PP^1)\setminus(B^4_+\sqcup B^4_-)$,
with two disjoint open balls. The exterior is
$Q\times S^1$, while each filling is
$S^3\times D^2$. Near an exceptional curve regard
a nonzero tautological vector as
$w\in\C^2\setminus0$. On the unit norm slice
$w\in S^3$, and the blowdown is
$u=a^m w$ or $u=(a')^m w$. The boundary maps are
\[
 \begin{split}
 (w,e^{\mathrm i\theta})&\longmapsto
       (e^{\mathrm i m\theta}w,e^{\mathrm i\theta}),\\
 (w,e^{\mathrm i\theta})&\longmapsto
       (e^{\mathrm i m\theta}w,e^{-\mathrm i\theta}).
 \end{split}                                             \tag{C.PE.18}
\]
The pair sequence for the two removed balls gives
free homology of $Q$ of ranks $(1,2,2,3,0)$; its
two boundary sphere classes are opposite
primitive classes $\delta,-\delta$ in $H_3(Q,\Z)$.
In Mayer--Vietoris for (C.PE.18), the map on
$H_1(S^3\times S^1)^{\oplus2}$ has primitive
rank-one image, the circle generator of the
exterior. The map on $H_3$ is injective with
primitive image, since each of the two filling
$S^3$ summands receives its own identity map.
The map on $H_4$ has primitive rank-one image
generated by $\delta\times[S^1]$. Every kernel
and cokernel used in the sequence is therefore
free, and it gives
\[
       \operatorname{rank}H_k(M_5,\Z)=(1,2,5,5,2,1).
                                                               \tag{C.PE.19}
\]
K\"unneth and (C.PE.16) give the $Y_m$ row of
(C.PE.3), including integral torsion-freeness.

Each Hopf surface is differentiably $S^3\times S^1$.
For the support use the actual union
$S_m=P_m\cup Q_+\cup Q_-$ with intersections
$J_+\sqcup J_-$. In the preceding $M_3$ presentation,
the core angular loop of $J_+$ is $\delta_+$ and
that of $J_-$ is $-\delta_+$. Their radial loops
are the $S^1$ factor of (C.PE.16). Inclusion in
$Q_\pm$ kills the angular loop and takes the
radial loop to its Hopf generator. Consequently
the integral map
\[
 \begin{split}
 H_1(J_+\sqcup J_-,\Z)\longrightarrow
 H_1(P_m,\Z)\oplus H_1(Q_+,\Z)\oplus H_1(Q_-,\Z)
 \end{split}                                             \tag{C.PE.20}
\]
has primitive rank three: its columns are, up to
orientation signs, $\delta_+$, $z-q_+$,
$-\delta_+$, $z-q_-$. On $H_2$, the two elliptic
fundamental classes map to the same primitive
fibre class of $P_m$ and to zero on the Hopf
surfaces. Primitivity follows directly from
$[J_\pm]=\pm\delta_+\times[S^1]$ with the appropriate
product orientations. More explicitly, the dual chart at
$J_-$ reverses both its angular loop and its radial loop:
its complex orientation is
$(-\delta_+)\times(-z)=\delta_+\times z$.
Using the same global radial generator $z$ at both
ends instead reverses the displayed product orientation
at $J_-$. The restriction map on
cohomology thus has rank one in degree two and
rank three in degree one. Mayer--Vietoris gives
the free support ranks $(1,3,6,7,3)$ and Euler zero.

The coherent groups also retain exact maps.
On every chart (C.PE.8), finite direct image and
the exact finite-group invariant functor give
\[
 \rho_*\OO_{Y_m}=\OO_V,\qquad
 R^1\rho_*\OO_{Y_m}=\OO_V,\qquad
 R^j\rho_*\OO_{Y_m}=0\quad(j\ge2).                       \tag{C.PE.21}
\]
Translations act trivially on $H^1(E,\OO_E)$.
The trivialization of $R^1$ is global: the
constant cocycle assigning $1$ to the generator
(C.PE.7) gives a class restricting to the
nonzero logarithmic cocycle on each elliptic
fibre, including the multiple fibres in the
finite-quotient chart. It is a nowhere-zero
frame. Since $h^k(V,\OO_V)=(1,1,0)$, Leray gives
$(1,2,1,0)$ for $Y_m$. The same local calculation
on $P_m\to E$ gives $(1,2,1)$ there, consistent
with (C.PE.2).

For a primary Hopf surface, a holomorphic
two-form lifts and extends to $\C^2$; invariance
under its strict scalar contraction forces
every power-series coefficient to vanish.
Hence $p_g=0$. Its topological $e=0$ and
$H^2(-,\Z)=0$ give $\chi(\OO)=0$ by Noether's
formula, so its coherent dimensions are
$(1,1,0)$. Restriction
$H^1(Q_\pm,\OO)\to H^1(J_\pm,\OO)$ is nonzero:
the deck cocycle $1$ restricts to the same
cocycle on $\C^*/q^{m\Z}$. It is not a
coboundary, since the constant Laurent
coefficient of $g(q^m u)-g(u)$ is zero for
every holomorphic $g$ on $\C^*$.
The coherent normalization sequence for
$S_m$ therefore has a surjective difference
map in degrees zero and one. Its remaining
terms give $h^k(S_m,\OO)=(1,2,1)$.

Finally, a rational curve in $P_m$ would lift
through the finite \'etale map (C.PE.12), since
$\PP^1$ is simply connected. Its two projections
to the positive-genus curves would be constant.
The Hopf components have no rational curve
either: a map from $\PP^1$ lifts to
$\C^2\setminus0$, and its coordinate functions
are constant. The inverse-map resolution
comparison in C.2.UN.3 prevents contraction
of any of these surfaces in a bimeromorphic
map between smooth threefold fillings. A
finite base change $t=s^d$ is generically,
along $P_m$, the graph $x=s^d$ for the local
transverse coordinate $x=t$. Its normalized
strict component maps with degree one to
$P_m$; resolution preserves that birational
component. The distinct plurigenus slopes
$2(1-1/m)$ distinguish the unique properly
elliptic components, and hence the fillings
for different $m$.

This is an explicit signed logarithmic-transform
example. The prior Kodaira-dimension-one
occurrences discussed after C.2.UN.6, including
\cite[Example~3.4(1)]{Fujimoto1988Log}, are retained
as the literature comparison; no priority claim
is made for this example.
\end{proof}
